\documentclass[a4paper,11pt]{article}
\usepackage[utf8]{inputenc}
\usepackage[T1]{fontenc}
\usepackage[dutch]{babel}
\usepackage{geometry}
\usepackage{listings}
\usepackage{xcolor}
\usepackage{hyperref}
\usepackage{graphicx}
\usepackage{placeins}
\usepackage{enumitem}
\usepackage{parskip}

\geometry{margin=2.5cm}

\definecolor{codegreen}{rgb}{0,0.6,0}
\definecolor{codegray}{rgb}{0.5,0.5,0.5}
\definecolor{codepurple}{rgb}{0.58,0,0.82}
\definecolor{backcolour}{rgb}{0.95,0.95,0.92}

\lstdefinestyle{mystyle}{
    backgroundcolor=\color{backcolour},
    commentstyle=\color{codegreen},
    keywordstyle=\color{magenta},
    numberstyle=\tiny\color{codegray},
    stringstyle=\color{codepurple},
    basicstyle=\ttfamily\footnotesize,
    breakatwhitespace=false,
    breaklines=true,
    captionpos=b,
    keepspaces=true,
    numbers=left,
    numbersep=5pt,
    showspaces=false,
    showstringspaces=false,
    showtabs=false,
    tabsize=2
}
\lstset{style=mystyle}

\title{Vorderingsstaten Agent\\[0.5em]\large Technische Beschrijving}
\author{AIDA Project Documentatie}
\date{\today}

\begin{document}
\maketitle
\tableofcontents
\newpage

%-----------------------------------------------------------------------
\section{Inleiding}
%-----------------------------------------------------------------------

Dit document beschrijft de werking van de \textbf{Vorderingsstaten Agent},
een softwaretool die automatisch vorderingsstaatrapporten genereert op
basis van werffoto's. De gebruiker uploadt twee foto's --- een situatie
\emph{voor} en een situatie \emph{na} --- en het systeem herkent welke
werken zijn uitgevoerd en koppelt deze aan de correcte posten in het
bouwtechnisch bestek. Het eindresultaat is een leesbaar Markdown-rapport
dat direct klaar is voor review.

Figuur~\ref{fig:multi-agent tool} toont een screenshot van de webinterface.

\begin{figure}[!htbp]
    \centering
    \includegraphics[width=0.5\linewidth]{Voorbeeld_tool.png}
    \caption{Screenshot van de multi-agent tool}
    \label{fig:multi-agent tool}
\end{figure}
\FloatBarrier

%-----------------------------------------------------------------------
\section{Kernconcepten}
%-----------------------------------------------------------------------

Lezers die niet vertrouwd zijn met de onderliggende technologieën vinden
hier een beknopte uitleg van de begrippen die in dit project een rol
spelen.

\subsection{Large Language Model (LLM) en Azure OpenAI}

Een \textbf{Large Language Model} is een AI-systeem dat getraind is op
enorme hoeveelheden tekst en daardoor complexe taal begrijpt en
genereert. Dit project maakt gebruik van model \texttt{gpt-5.2}, gehost
op het Microsoft Azure-platform. Een groot voordeel van dit specifieke
model is dat het ook \emph{afbeeldingen} kan verwerken: het ``kijkt''
naar een werffoto en beschrijft wat er zichtbaar is.

Azure OpenAI is de Microsoft-variant van OpenAI's API. Het gedrag is
nagenoeg identiek aan de standaard OpenAI-dienst, maar de verbinding
loopt via een eigen bedrijfsendpoint. Dat is relevant voor
beveiligingsbeleid en factuurbeheer.

\subsection{Vector Store en semantisch zoeken}

Een bouwtechnisch bestek beslaat al snel honderden pagina's. Een AI-model
kan niet al die tekst tegelijk in zijn ``geheugen'' houden. De oplossing
is een \textbf{Vector Store}: een soort intelligente index die het bestek
in kleine fragmenten opslaat.

Elk fragment wordt omgezet in een reeks getallen --- een \emph{vector}
--- die de \emph{betekenis} van de tekst vastlegt. Wanneer het model wil
opzoeken welke bestekeisen gelden voor ``dakisolatie'', zoekt de Vector
Store niet op het exacte woord, maar op fragmenten die \emph{inhoudelijk
verwant} zijn. Dit heet \textbf{semantisch zoeken} en levert veel
relevantere resultaten op dan een gewone trefwoordzoekopdracht.

In dit project bevatten de Vector Store twee soorten documenten:
\begin{itemize}
    \item De \textbf{master-inhoudstafel} met alle bestekpostnummers en hun
          omschrijving --- als snelle index.
    \item De \textbf{tien deelbestekken} (Deel-0 t/m Deel-9) met de volledige
          technische eisen per post.
\end{itemize}

\subsection{Retrieval-Augmented Generation (RAG)}

RAG is de techniek die hallucinaties (het verzinnen van feiten) sterk
vermindert. In plaats van puur op zijn getrainde kennis te steunen,
\emph{zoekt} het model eerst de relevante bestekfragmenten op voordat het
een antwoord genereert. Het rapport is daardoor gebaseerd op de werkelijke
bestektekst, niet op aannames.

%-----------------------------------------------------------------------
\section{Werking van de applicatie}
%-----------------------------------------------------------------------

Het systeem doorloopt vijf opeenvolgende stappen voor elke aanvraag.

\subsection{Stap 1 --- Input en optimalisatie}

De gebruiker uploadt twee foto's via een eenvoudige webpagina of via de
API. Omdat werffoto's van een moderne smartphone al snel 10--20~MB groot
zijn, worden ze automatisch verkleind (maximaal 2048$\times$2048~pixels)
en gecomprimeerd. Dit versnelt de verwerking en verlaagt de kosten.

Vóór de foto's worden verwerkt controleert het systeem ook of er
voldoende geheugen en rekenkracht beschikbaar zijn. Dreigt de server
overbelast te raken, dan wordt de aanvraag netjes geweigerd met een
foutmelding in plaats van onverwacht te crashen.

\subsection{Stap 2 --- Agent 1: de visuele inspecteur}

De eerste AI-agent bekijkt de twee foto's en analyseert het visuele
verschil. Zijn enige taak is observeren en classificeren. Hij zoekt in de
master-inhoudstafel naar de \emph{bestekpostnummers} die het beste
aansluiten bij wat er zichtbaar is bijgekomen, verwijderd of gewijzigd.

De output van Agent~1 is een strikt gestructureerd JSON-document met:
\begin{itemize}
    \item de gevonden bestekpostnummers (bijv.\ \texttt{30.14}),
    \item korte visuele observaties per overgang,
    \item een zekerheidsscore: \emph{hoog}, \emph{middel} of \emph{laag}.
\end{itemize}

Agent~1 mag uitsluitend nummers gebruiken die letterlijk in de
master-inhoudstafel staan. Twijfelgevallen worden weggelaten of
teruggebracht tot een meer algemene post.

\subsection{Stap 3 --- Agent 2: de rapporteur}

De tweede AI-agent ontvangt de output van Agent~1 en gaat aan de slag
met het detailbestek. Via de \texttt{file\_search}-tool raadpleegt hij
automatisch de relevante deelbestekken om de technische eisen op te
halen die bij elk gevonden bestekpostnummer horen.

Op basis van die informatie schrijft Agent~2 het definitieve
vorderingsstaatrapport. Per bestekpost beschrijft hij wat er zichtbaar
is uitgevoerd, welke bestekeisen van toepassing zijn, wat de bronreferentie
is in het bestek, en welke open punten of risico's er bestaan.

\subsection{Stap 4 --- MVP-regel bij lage zekerheid}

Na het rapport van Agent~2 controleert het systeem de zekerheidsscore van
Agent~1. Is die \emph{laag}, dan wordt automatisch de melding
``\textbf{Manuele check / extra foto nodig}'' toegevoegd aan de relevante
bestekposten. Dit is een harde regel die niet door de AI overgeslagen kan
worden.

\subsection{Stap 5 --- Output en logging}

Het eindrapport wordt teruggegeven als leesbaar Markdown-document en is
direct te bekijken in de webinterface. Tegelijk slaat het systeem een
gedetailleerd logbestand op (in \texttt{.jsonl}-formaat) met de volledige
tussentijdse output van beide agenten, alle tijdsmetingen per stap, en
een unieke aanvraag-ID. Dat maakt het achteraf eenvoudig om de kwaliteit
van de output te beoordelen en problemen te traceren.

%-----------------------------------------------------------------------
\section{Eenmalige initialisatie: kennisbank opbouwen}
%-----------------------------------------------------------------------

Voordat de applicatie gebruikt kan worden, moet de Vector Store eenmalig
worden gevuld. Dit gebeurt via het script \texttt{scripts/init\_kb.py}.
Het script:
\begin{enumerate}
    \item maakt één Vector Store aan op Azure,
    \item uploadt de master-inhoudstafel met de juiste metadata,
    \item uploadt alle tien deelbestekken elk met hun eigen metadatalabel
          (\texttt{deel=0} t/m \texttt{deel=9}),
    \item wacht tot elk bestand volledig geïndexeerd is, en
    \item schrijft de unieke \texttt{VECTOR\_STORE\_ID} weg naar het
          \texttt{.env}-configuratiebestand.
\end{enumerate}

Dit hoeft maar één keer uitgevoerd te worden, tenzij de bestekken worden
bijgewerkt.

%-----------------------------------------------------------------------
\section{Systeeminstructies}
%-----------------------------------------------------------------------

De intelligentie van de agenten wordt aangestuurd via zogenaamde
\emph{system prompts}: een set instructies die de AI bij elke aanvraag
meegeeft. Hieronder staan de exacte instructies voor beide agenten.

\subsection{Agent 1 --- Visie en bestekpostmapping}

\lstinputlisting[
  language={},
  caption={Systeeminstructies Agent 1 (prompts/agent1\_system.txt)},
  label={lst:agent1}
]{prompts/agent1_system.txt}

\subsection{Agent 2 --- Rapportage}

\lstinputlisting[
  language={},
  caption={Systeeminstructies Agent 2 (prompts/agent2\_system.txt)},
  label={lst:agent2}
]{prompts/agent2_system.txt}

%-----------------------------------------------------------------------
\section{Configuratie}
%-----------------------------------------------------------------------

De applicatie heeft slechts twee configuratiewaarden nodig die de
gebruiker zelf moet invullen in een \texttt{.env}-bestand:
\begin{itemize}
    \item \texttt{AZURE\_OPENAI\_API\_KEY} --- de persoonlijke sleutel
          voor de Azure-dienst.
    \item \texttt{VECTOR\_STORE\_ID} --- wordt automatisch ingevuld door
          het initialisatiescript.
\end{itemize}

Alle overige instellingen (endpoint-URL, deployment-naam, API-versie)
zijn vastgelegd in de broncode en hoeven niet apart geconfigureerd te
worden.

%-----------------------------------------------------------------------
\section{Conclusie}
%-----------------------------------------------------------------------

Door de taken te verdelen over twee gespecialiseerde agenten --- één die
kijkt, en één die schrijft --- combineert het systeem visuele analyse met
diepgaande bestek\-kennis. De Vector Store garandeert dat het rapport
gebaseerd is op de werkelijke bestektekst. De ingebouwde
zekerheidsbeoordeling en verplichte waarschuwingen bij lage zekerheid
zorgen ervoor dat de output bruikbaar en eerlijk is, ook wanneer de
foto-invoer onvolledig is.

\end{document}
